% ============================================================
%  LEMBAR SOAL  --  FISIKA  --  SSU-ITB 2024  --  Paket 1
%  Tipe A (1-20) + Tipe B (21-39) + Soal Baru (40)
%  Komposisi: 40% PG Biasa | 40% PG Majemuk | 20% Isian
% ============================================================

\documentclass[11pt]{article}

\usepackage[utf8]{inputenc}
\usepackage[T1]{fontenc}
\usepackage[a4paper,margin=2.5cm,top=2.8cm]{geometry}
\usepackage{amsmath,amssymb}
\usepackage{graphicx}
\usepackage{enumitem}
\usepackage{multicol}
\usepackage{xcolor}
\usepackage[most]{tcolorbox}
\usepackage{fancyhdr}
\usepackage{booktabs}
\usepackage{icomma}
\usepackage{array}

\pagestyle{fancy}\fancyhf{}
\lhead{\small SSU-ITB 2024 --- Fisika}
\rhead{\small Paket 1}
\cfoot{\small\thepage}
\renewcommand{\headrulewidth}{0.4pt}

\newtcolorbox{soal}[1]{
  breakable,
  colback=white, colframe=black, colbacktitle=black,
  coltitle=white, fonttitle=\bfseries\sffamily,
  title=Nomor #1,
  boxrule=0.8pt, arc=0pt,
  left=3mm, right=3mm, top=2mm, bottom=2mm}

\newenvironment{pil}{%
  \begin{enumerate}[label=(\Alph*),leftmargin=*,itemsep=1pt,topsep=3pt]}%
  {\end{enumerate}}
\newenvironment{pilB}{%
  \par\smallskip\begin{multicols}{2}
  \begin{enumerate}[label=(\Alph*),leftmargin=*,itemsep=1pt,topsep=3pt]}%
  {\end{enumerate}\end{multicols}}

\newcommand{\gambar}[2][0.52\textwidth]{%
  \begin{center}\includegraphics[width=#1]{#2}\end{center}}

\linespread{1.05}

% ===================================================================
\begin{document}
\thispagestyle{empty}

\begin{center}
  \fbox{\parbox{0.6\textwidth}{\centering\bfseries\large LEMBAR SOAL}}
\end{center}
\vspace{0.5cm}

\begin{center}
  {\LARGE\bfseries FISIKA}
\end{center}
\vspace{0.5cm}

\begin{center}
  \begin{tabular}{@{\hspace{2cm}}ll@{\hspace{0.4cm}:@{\hspace{0.4cm}}}l}
    Jenjang      & & SMA / Sederajat               \\[0.3em]
    Hari/Tanggal & & \underline{\hspace{5cm}}      \\[0.3em]
    Waktu        & & \underline{\hspace{5cm}}      \\[0.3em]
    Paket Soal   & & 1                             \\
  \end{tabular}
\end{center}

\vspace{0.6cm}
\noindent\rule{\linewidth}{0.6pt}
\vspace{0.3cm}

\noindent\textbf{\underline{PETUNJUK UMUM}}
\begin{enumerate}[leftmargin=*,itemsep=4pt,topsep=4pt]
  \item Anda \textbf{tidak} diperbolehkan menggunakan sumber-sumber eksternal,
        termasuk internet, \textit{large language model} seperti ChatGPT, Claude,
        LLaMA, Gemini, Meta AI, dan sebagainya. Anda sangat dianjurkan untuk
        mencoba mengerjakan sendiri terlebih dahulu karena evaluasi ini dibuat
        untuk \textbf{mengukur} tingkat kepahaman; nilai $<70$ mengindikasikan
        \textbf{tidak lulus}, dan jika sudah melewati \textit{deadline} maka
        dianggap tidak mengerjakan yang berarti \textbf{tidak lulus}.

  \item Terdapat 2 tahap pengerjaan, yaitu tahap \textbf{Try Out} dan tahap
        \textbf{Review}. Try Out bersifat \textit{open book}; Anda diperbolehkan
        membuka buku apapun selama bersifat \textit{offline} dengan rentang waktu
        2 jam pada tanggal \underline{\hspace{3cm}}. Setelah itu memasuki
        tahap Review dengan \textit{schedule} rentang tanggal
        \underline{\hspace{3cm}}--\underline{\hspace{3cm}}, di mana Anda
        diharapkan mengerjakan ulang dengan disertakan penjelasan/pembelaan
        mengapa revisi jawaban adalah pilihan tersebut, dan tetap mencantumkan
        penjelasan pada soal yang walaupun pada penilaian sudah benar.

  \item Anda dianjurkan untuk mencantumkan caranya walaupun tidak termasuk
        penilaian, mengingat sifatnya adalah sebagai bahan pembelajaran.
        Mohon bertanggung jawab atas pekerjaan Anda sendiri.

  \item Bacalah setiap soal dengan teliti. Terdapat tiga jenis soal:
        \begin{enumerate}[label=(\arabic*),leftmargin=2em,topsep=2pt,itemsep=2pt]
          \item \textbf{Pilihan Ganda Biasa} --- 5 opsi (A--E), pilih
                1 jawaban paling tepat;
          \item \textbf{Pilihan Ganda Majemuk Kompleks} --- tabel
                \textbf{Benar}/\textbf{Salah}, tentukan kebenaran tiap
                pernyataan secara mandiri; dan
          \item \textbf{Isian Singkat} --- tulis nilai/ekspresi pada
                kolom yang tersedia.
        \end{enumerate}

  \item Soal berjumlah \textbf{40 butir} soal dengan bobot asli per soal adalah
        \textbf{2{,}5} untuk mendapatkan nilai sempurna sebesar \textbf{100}.
        Gunakan $g=10$ m/s$^2$ bila tidak disebutkan lain.

  \item Silakan bertanya apabila terdapat soal yang ambigu, rancu, atau
        kurang spesifik selama itu tidak menjurus kepada jawaban soal
        tersebut.

  \item Isilah detail informasi berikut.\\[0.3em]
        \textbf{Nama}\hspace{0.45cm};\enspace\underline{\hspace{8cm}}
\end{enumerate}

\vspace{0.6cm}
\noindent\textit{Dengan menandatangani kolom di bawah ini, saya menyatakan telah membaca,
memahami, dan menyetujui seluruh peraturan yang tercantum di atas.}

\vspace{0.6cm}
\noindent\fbox{\parbox{5cm}{\vspace{2.2cm}\centering\small Tanda Tangan Peserta\vspace{0.3cm}}}

\clearpage

% ===================================================================
%  TIPE A  (nomor 1--20)
% ===================================================================

\begin{soal}{1}  % PG Majemuk
Sebuah partikel bergerak sepanjang sumbu $x$. Pada $t=0$, partikel
mulai bergerak ke arah sumbu $x$ positif. Grafik kecepatan terhadap
waktu partikel tersebut ditunjukkan berikut.
\gambar[0.62\textwidth]{images/1/1-1.png}
Tentukan kebenaran setiap pernyataan berikut!

\par\smallskip
{\renewcommand{\arraystretch}{1.4}%
\begin{tabular}{@{}p{0.76\linewidth}cc@{}}
\hline
\textbf{Pernyataan} & \textbf{Benar} & \textbf{Salah}\\
\hline
(1)~Pada selang waktu 20--60\,s, jarak tempuh partikel sama dengan nol.
  & $\bigcirc$ & $\bigcirc$\\
(2)~Pada selang waktu 20--60\,s, percepatan partikel tetap (konstan).
  & $\bigcirc$ & $\bigcirc$\\
(3)~Pada selang waktu 40--60\,s, partikel bergerak diperlambat.
  & $\bigcirc$ & $\bigcirc$\\
(4)~Kelajuan partikel pada selang 10--20\,s sama dengan kelajuannya
    pada selang 60--70\,s.
  & $\bigcirc$ & $\bigcirc$\\
\hline
\end{tabular}}
\end{soal}

\begin{soal}{2}  % PG Biasa
Sebuah partikel bermassa 500 gram mula-mula bergerak dengan kelajuan
2,8 m/s. Partikel tersebut kemudian ditarik oleh gaya konstan $F$
sehingga dipercepat 2 m/s$^2$ selama 8 detik. Besar usaha yang
dilakukan oleh gaya $F$ selama dipercepat adalah \ldots\ J.
\begin{pilB}
  \item $72$ \item $76{,}8$ \item $80$ \item $86{,}4$ \item $96$
\end{pilB}
\end{soal}

\begin{soal}{3}  % PG Biasa
Sebuah bola dilepaskan dari atas gedung dan pada saat yang bersamaan
sebuah batu dilemparkan secara horizontal dari ketinggian yang sama.
Jika gesekan udara diabaikan, pernyataan yang paling benar adalah \ldots
\begin{pil}
  \item Bola mencapai tanah lebih dahulu karena tidak memiliki
        kecepatan horizontal.
  \item Batu mencapai tanah lebih dahulu karena memiliki tambahan
        kecepatan horizontal.
  \item Keduanya mencapai tanah pada saat yang sama.
  \item Laju akhir bola dan batu sesaat sebelum menyentuh tanah
        adalah sama.
  \item Trajektori (lintasan) bola dan batu adalah identik.
\end{pil}
\end{soal}

\begin{soal}{4}  % PG Majemuk
Pada sebuah eksperimen diperoleh data berupa jari-jari lintasan dan
percepatan sentripetal dari sebuah partikel yang bergerak melingkar
beraturan.
\begin{center}
\begin{tabular}{ccc}
\hline
No & Jari-jari (m) & $a_s$ (m/s$^2$)\\
\hline
1 & 0,2 & 45 \\
2 & 0,4 & 22,5 \\
3 & 0,6 & 15 \\
4 & 0,8 & 11,25 \\
5 & 1   & 9 \\
\hline
\end{tabular}
\end{center}
Tentukan kebenaran setiap pernyataan berikut!

\par\smallskip
{\renewcommand{\arraystretch}{1.4}%
\begin{tabular}{@{}p{0.76\linewidth}cc@{}}
\hline
\textbf{Pernyataan} & \textbf{Benar} & \textbf{Salah}\\
\hline
(1)~Hubungan $v^2 = a_s \cdot r$ berlaku pada setiap baris data
    dalam tabel.
  & $\bigcirc$ & $\bigcirc$\\
(2)~Kelajuan linear partikel pada eksperimen tersebut adalah 3 m/s.
  & $\bigcirc$ & $\bigcirc$\\
(3)~Percepatan sentripetal berbanding lurus dengan jari-jari lintasan
    saat kelajuan partikel konstan.
  & $\bigcirc$ & $\bigcirc$\\
(4)~Kelajuan sudut $\omega$ partikel sama pada setiap baris data.
  & $\bigcirc$ & $\bigcirc$\\
\hline
\end{tabular}}
\end{soal}

\begin{soal}{5}  % PG Biasa
Sebuah benda bergerak melingkar beraturan. Pernyataan yang benar
mengenai percepatan benda adalah \ldots
\begin{pil}
  \item Percepatannya nol karena kelajuannya tetap.
  \item Percepatan selalu searah dengan kecepatan benda.
  \item Percepatan selalu menuju pusat lingkaran, besarnya tetap,
        tetapi arahnya berubah.
  \item Percepatan selalu menuju pusat lingkaran dan besarnya berubah.
  \item Percepatan menjauh dari pusat lingkaran dengan besar tetap.
\end{pil}
\end{soal}

\begin{soal}{6}  % PG Biasa
Dua buah balok bermassa $m_1=2$ kg dan $m_2=3$ kg terikat pada tali
yang melalui sebuah katrol. Massa katrol, massa tali, dan gesekan
katrol diabaikan. Balok $m_1$ berada di atas meja horizontal kasar
dengan $\mu_s=0{,}5$ dan $\mu_k=0{,}4$.
\gambar[0.42\textwidth]{images/1/1-2.png}
Besar gaya tegangan tali adalah \ldots\ N.
\begin{pilB}
  \item $8{,}0$ \item $12{,}4$ \item $14{,}0$ \item $16{,}8$ \item $22{,}4$
\end{pilB}
\end{soal}

\begin{soal}{7}  % PG Majemuk
Sebuah balok bermassa 2,5 kg meluncur turun sejauh 1 meter dari
puncak bidang miring dengan sudut kemiringan $\alpha$
$\left(\tan\alpha=\tfrac{3}{4}\right)$ terhadap horizontal.
Koefisien gesek statis dan kinetis masing-masing 0,5 dan 0,2.
Tentukan kebenaran setiap pernyataan berikut!

\par\smallskip
{\renewcommand{\arraystretch}{1.4}%
\begin{tabular}{@{}p{0.76\linewidth}cc@{}}
\hline
\textbf{Pernyataan} & \textbf{Benar} & \textbf{Salah}\\
\hline
(1)~Komponen gaya gravitasi sepanjang bidang miring adalah
    $mg\sin\alpha=15$ N.
  & $\bigcirc$ & $\bigcirc$\\
(2)~Gaya gesek kinetik pada bidang miring adalah $f_k=4$ N.
  & $\bigcirc$ & $\bigcirc$\\
(3)~Usaha yang dilakukan gaya gesek bernilai positif.
  & $\bigcirc$ & $\bigcirc$\\
(4)~Usaha gaya nonkonservatif selama perpindahan 1 m adalah $-4$ J.
  & $\bigcirc$ & $\bigcirc$\\
\hline
\end{tabular}}
\end{soal}

\begin{soal}{8}  % PG Biasa
Seorang anak laki-laki dengan berat 50 kg berdiri di dalam lift.
\gambar[0.32\textwidth]{images/1/1-3.png}
Perhatikan pernyataan-pernyataan tentang gaya yang bekerja pada kaki
anak tersebut:
\begin{enumerate}[label=\alph*.,leftmargin=*,itemsep=3pt,topsep=3pt]
  \item Jika lift diam, gaya pada kaki anak sama dengan beratnya.
  \item Jika lift dipercepat ke atas, gaya pada kaki anak lebih besar
        dari beratnya dan arah gayanya ke atas.
  \item Jika lift bergerak ke atas dengan kecepatan tetap, gaya pada
        kaki anak lebih kecil dari beratnya dan arahnya ke atas.
  \item Jika lift dipercepat ke bawah, gaya pada kaki anak lebih kecil
        dari beratnya dan arahnya ke bawah.
\end{enumerate}
Pernyataan yang benar adalah \ldots
\begin{pil}
  \item a, b, dan c
  \item a, b, c, dan d
  \item a saja
  \item a dan b
  \item Tidak ada pernyataan yang benar
\end{pil}
\end{soal}

\begin{soal}{9}  % PG Majemuk
Dua buah partikel identik A dan B masing-masing sedang bergerak
melingkar. Nilai percepatan sentripetal terhadap jari-jari lintasan
dari kedua partikel ditunjukkan pada grafik berikut.
\gambar[0.42\textwidth]{images/1/1-4.png}
Tentukan kebenaran setiap pernyataan berikut!

\par\smallskip
{\renewcommand{\arraystretch}{1.4}%
\begin{tabular}{@{}p{0.76\linewidth}cc@{}}
\hline
\textbf{Pernyataan} & \textbf{Benar} & \textbf{Salah}\\
\hline
(1)~Kemiringan garis A pada grafik $a_s$ vs $r$ menyatakan
    $\omega_A^2=1$ (rad/s)$^2$.
  & $\bigcirc$ & $\bigcirc$\\
(2)~Kemiringan garis B pada grafik $a_s$ vs $r$ menyatakan
    $\omega_B^2=2$ (rad/s)$^2$.
  & $\bigcirc$ & $\bigcirc$\\
(3)~Kedua partikel memiliki kelajuan sudut yang sama,
    $\omega_A=\omega_B=1$ rad/s.
  & $\bigcirc$ & $\bigcirc$\\
(4)~Selisih kelajuan sudut kedua partikel adalah $(\sqrt{2}-1)$ rad/s.
  & $\bigcirc$ & $\bigcirc$\\
\hline
\end{tabular}}
\end{soal}

\begin{soal}{10}  % Isian
Sebuah balok yang diam bermassa 4,7 kg ditarik oleh gaya horizontal
konstan $F=10$ N di atas permukaan horizontal kasar dengan
$\mu_s=0{,}6$ dan $\mu_k=0{,}4$.
\gambar[0.5\textwidth]{images/1/1-5.png}
Percepatan balok tersebut adalah \ldots\ m/s$^2$.
\end{soal}

\begin{soal}{11}  % PG Majemuk
Perhatikan pernyataan-pernyataan berikut tentang momentum dan gerak
pusat massa.
Tentukan kebenaran setiap pernyataan berikut!

\par\smallskip
{\renewcommand{\arraystretch}{1.4}%
\begin{tabular}{@{}p{0.76\linewidth}cc@{}}
\hline
\textbf{Pernyataan} & \textbf{Benar} & \textbf{Salah}\\
\hline
(1)~Momentum sistem bisa kekal meskipun energi mekanik sistem
    tidak kekal.
  & $\bigcirc$ & $\bigcirc$\\
(2)~Syarat kekekalan momentum sistem adalah resultan gaya luar
    pada sistem bernilai nol.
  & $\bigcirc$ & $\bigcirc$\\
(3)~Kecepatan pusat massa sistem berubah hanya ketika ada resultan
    gaya luar yang bekerja.
  & $\bigcirc$ & $\bigcirc$\\
(4)~Momentum total sistem selalu kekal dalam segala proses.
  & $\bigcirc$ & $\bigcirc$\\
\hline
\end{tabular}}
\end{soal}

\begin{soal}{12}  % PG Majemuk
Dua bola identik A dan B dilepaskan dari atap gedung. Bola A
dilepaskan 1 detik sebelum bola B. Hambatan udara diabaikan dan
$g=10$ m/s$^2$.
Tentukan kebenaran setiap pernyataan berikut (ditinjau saat bola B
baru saja dilepaskan)!

\par\smallskip
{\renewcommand{\arraystretch}{1.4}%
\begin{tabular}{@{}p{0.76\linewidth}cc@{}}
\hline
\textbf{Pernyataan} & \textbf{Benar} & \textbf{Salah}\\
\hline
(1)~Kecepatan bola A saat B dilepaskan adalah 10 m/s ke bawah.
  & $\bigcirc$ & $\bigcirc$\\
(2)~Bola A sudah berada 5 m di bawah titik awal saat B dilepaskan.
  & $\bigcirc$ & $\bigcirc$\\
(3)~Setelah B dilepaskan, kecepatan relatif bola A terhadap bola B
    berkurang seiring waktu.
  & $\bigcirc$ & $\bigcirc$\\
(4)~Setelah B dilepaskan, jarak vertikal antara kedua bola selalu
    tetap 5 m.
  & $\bigcirc$ & $\bigcirc$\\
\hline
\end{tabular}}
\end{soal}

\begin{soal}{13}  % PG Majemuk
Sebuah balok bermassa 2,5 kg meluncur dari puncak bidang miring
$\left(\tan\alpha=\tfrac{3}{4}\right)$ dengan laju awal 1 m/s.
Koefisien gesek statis dan kinetis masing-masing 0,50 dan 0,20.
Tentukan kebenaran setiap pernyataan berikut!

\par\smallskip
{\renewcommand{\arraystretch}{1.4}%
\begin{tabular}{@{}p{0.76\linewidth}cc@{}}
\hline
\textbf{Pernyataan} & \textbf{Benar} & \textbf{Salah}\\
\hline
(1)~Percepatan balok sepanjang bidang miring adalah $4{,}4$ m/s$^2$.
  & $\bigcirc$ & $\bigcirc$\\
(2)~Kecepatan awal balok di puncak bidang miring adalah 1 m/s.
  & $\bigcirc$ & $\bigcirc$\\
(3)~Setelah meluncur sejauh 1 m, kelajuan balok adalah 4,8 m/s.
  & $\bigcirc$ & $\bigcirc$\\
(4)~Nilai $v^2$ setelah meluncur sejauh 1 m adalah 9,8 m$^2$/s$^2$.
  & $\bigcirc$ & $\bigcirc$\\
\hline
\end{tabular}}
\end{soal}

\begin{soal}{14}  % PG Biasa
Dua buah balok dengan massa masing-masing $m=7{,}6$ kg dan $M=2$ kg
dihubungkan oleh tali yang melalui sebuah katrol. Koefisien gesek
statis antara dua permukaan balok adalah $\mu_s=1{,}1$, sedangkan
lantai licin.
\gambar[0.5\textwidth]{images/1/1-6.png}
Besar gaya horizontal maksimum $F$ sehingga sistem masih dalam
kesetimbangan adalah \ldots\ N.
\begin{pilB}
  \item $83{,}6$ \item $120{,}0$ \item $136{,}0$ \item $167{,}2$ \item $200{,}0$
\end{pilB}
\end{soal}

\begin{soal}{15}  % PG Biasa
Dadang melakukan perjalanan dari kota A ke kota B dengan kecepatan
rata-rata 58 km/jam, kemudian kembali dari kota B ke kota A dengan
kecepatan rata-rata 75 km/jam. Laju rata-rata seluruh perjalanan
Dadang adalah \ldots\ km/jam.
\begin{pilB}
  \item $63{,}5$ \item $65{,}0$ \item $65{,}4$ \item $66{,}5$ \item $67{,}5$
\end{pilB}
\end{soal}

\begin{soal}{16}  % PG Majemuk
Grafik di bawah menyatakan simpangan gelombang berjalan pada $t=0$
yang merambat ke kanan. Frekuensi gelombang adalah 5 Hz.
\gambar[0.55\textwidth]{images/1/1-7.png}
Tentukan kebenaran setiap pernyataan berikut!

\par\smallskip
{\renewcommand{\arraystretch}{1.4}%
\begin{tabular}{@{}p{0.76\linewidth}cc@{}}
\hline
\textbf{Pernyataan} & \textbf{Benar} & \textbf{Salah}\\
\hline
(1)~Amplitudo gelombang tersebut adalah 2 m.
  & $\bigcirc$ & $\bigcirc$\\
(2)~Panjang gelombang $\lambda = 2$ m.
  & $\bigcirc$ & $\bigcirc$\\
(3)~Bilangan gelombang $k = 2\pi$ rad/m.
  & $\bigcirc$ & $\bigcirc$\\
(4)~Persamaan simpangan gelombang adalah
    $y=2\cos(\pi x - 10\pi t)$.
  & $\bigcirc$ & $\bigcirc$\\
\hline
\end{tabular}}
\end{soal}

\begin{soal}{17}  % Isian
Gelombang bunyi dengan frekuensi 200 Hz dimasukkan ke bagian atas
tabung vertikal berisi air. Gelombang berdiri dihasilkan pada dua
ketinggian air berturut-turut 30 cm dan 90 cm. Laju gelombang bunyi
pada kolom udara di tabung tersebut adalah \ldots\ m/s.
\end{soal}

\begin{soal}{18}  % PG Majemuk
Sebuah gelombang berjalan merambat sepanjang tali dengan persamaan
simpangan
\[
  y(x,t) = (2\,\text{m})\sin\!\left(2\pi x - 3\pi t - \frac{\pi}{3}\right)
\]
dengan $t$ dalam sekon dan $x$ dalam meter.
Tentukan kebenaran setiap pernyataan berikut!

\par\smallskip
{\renewcommand{\arraystretch}{1.4}%
\begin{tabular}{@{}p{0.76\linewidth}cc@{}}
\hline
\textbf{Pernyataan} & \textbf{Benar} & \textbf{Salah}\\
\hline
(1)~Gelombang merambat ke arah sumbu $+x$.
  & $\bigcirc$ & $\bigcirc$\\
(2)~Amplitudo gelombang adalah 2 m.
  & $\bigcirc$ & $\bigcirc$\\
(3)~Frekuensi sudut gelombang adalah $\omega = 2\pi$ rad/s.
  & $\bigcirc$ & $\bigcirc$\\
(4)~Laju rambat gelombang adalah 1,5 m/s.
  & $\bigcirc$ & $\bigcirc$\\
\hline
\end{tabular}}
\end{soal}

\begin{soal}{19}  % Isian
Persamaan simpangan sebuah gelombang berjalan transversal pada tali
adalah
\[
  y(x,t) = 0{,}28\sin\!\left(\frac{\pi x}{5} + 8\pi t\right)
\]
dengan $t$ dalam sekon dan $x$ dalam meter. Laju transversal tali di
posisi $x=1{,}5$ m saat $t=0{,}5$ s adalah \ldots\ m/s.
\end{soal}

\begin{soal}{20}  % PG Biasa  [SOAL BARU: Doppler mendekati dinding]
Sebuah sumber bunyi dengan frekuensi 500 Hz bergerak mendekati
pengamat yang diam dengan kecepatan 20 m/s. Cepat rambat bunyi di
udara adalah 340 m/s. Frekuensi bunyi yang terdeteksi oleh pengamat
adalah \ldots\ Hz.
\begin{pilB}
  \item $468{,}75$ \item $500{,}00$ \item $512{,}50$ \item $531{,}25$ \item $562{,}50$
\end{pilB}
\end{soal}

% ===================================================================
%  TIPE B  (nomor 21--40)
% ===================================================================

\begin{soal}{21}  % PG Biasa
Seorang pelempar akan melemparkan sebuah bola dengan laju 30 m/s
dan sudut elevasi $\alpha$ $\left(\tan\alpha=\tfrac{3}{4}\right)$.
Seorang penangkap berada pada jarak 60 m dari pelempar. Agar
penangkap dapat menangkap bola, dia harus berlari dengan laju
minimum sebesar \ldots\ m/s.
\begin{pilB}
  \item $5{,}50$ \item $6{,}00$ \item $6{,}50$ \item $7{,}33$ \item $8{,}00$
\end{pilB}
\end{soal}

\begin{soal}{22}  % PG Biasa
Dua buah balok A ($m_A=2{,}7$ kg) dan B ($m_B=5{,}9$ kg) terhubung
oleh tali. Balok A terhubung pada pasak melalui tali sepanjang 8,2 m;
balok B terhubung pada balok A melalui tali 5,1 m. Sistem bergerak
melingkar pada meja licin dengan $\omega=15$ rad/s.
\gambar[0.5\textwidth]{images/1/1-8.png}
Perbandingan $T_a / T_b$ adalah \ldots
\begin{pilB}
  \item $1{,}00$ \item $1{,}20$ \item $1{,}28$ \item $1{,}50$ \item $1{,}75$
\end{pilB}
\end{soal}

\begin{soal}{23}  % PG Biasa
Sebuah partikel bermassa 4 gram bertumbukan dengan partikel lain
bermassa 2 gram yang diam. Keduanya menempel setelah tumbukan.
Perbandingan energi yang hilang terhadap energi kinetik awal adalah
\ldots
\begin{pilB}
  \item $\dfrac{1}{6}$ \item $\dfrac{1}{4}$ \item $\dfrac{1}{3}$
  \item $\dfrac{1}{2}$ \item $\dfrac{2}{3}$
\end{pilB}
\end{soal}

\begin{soal}{24}  % PG Majemuk
Pada $t=0$ sebuah partikel diam di $x=0$ m. Partikel bergerak dengan
percepatan seperti pada grafik berikut.
\gambar[0.55\textwidth]{images/1/1-9.png}
Tentukan kebenaran setiap pernyataan berikut!

\par\smallskip
{\renewcommand{\arraystretch}{1.4}%
\begin{tabular}{@{}p{0.76\linewidth}cc@{}}
\hline
\textbf{Pernyataan} & \textbf{Benar} & \textbf{Salah}\\
\hline
(1)~Pada selang waktu 0--5\,s, partikel bergerak ke arah sumbu
    $x$ positif.
  & $\bigcirc$ & $\bigcirc$\\
(2)~Pada $t=5$ s dan $t=10$ s, partikel memiliki kecepatan yang sama.
  & $\bigcirc$ & $\bigcirc$\\
(3)~Pada selang waktu 5--10\,s, partikel diam (tidak bergerak).
  & $\bigcirc$ & $\bigcirc$\\
(4)~Pada $t=12$ s, partikel berada di posisi sumbu $x$ negatif.
  & $\bigcirc$ & $\bigcirc$\\
\hline
\end{tabular}}
\end{soal}

\begin{soal}{25}  % Isian
Sebuah partikel bergerak dari pusat koordinat ke arah sumbu $x$
positif dengan kecepatan 4 m/s selama 50 sekon, kemudian 5 m/s
selama 20 sekon berikutnya, lalu berbalik arah dengan kecepatan
3 m/s selama 30 sekon. Besar kecepatan rata-rata partikel dalam
100 sekon tersebut adalah \ldots\ m/s.
\end{soal}

\begin{soal}{26}  % PG Biasa
Balok bermassa 5 kg ditekan pada dinding dengan gaya horizontal
$F$. Koefisien gesek antara balok dan dinding adalah 0,3.
\gambar[0.5\textwidth]{images/1/1-10.png}
Gaya minimum yang diperlukan agar balok tidak jatuh adalah \ldots\ N.
\begin{pilB}
  \item $100{,}0$ \item $125{,}0$ \item $133{,}3$ \item $166{,}7$ \item $200{,}0$
\end{pilB}
\end{soal}

\begin{soal}{27}  % PG Majemuk
Suatu bola dijatuhkan secara tegak lurus dari ketinggian 50 m dan
kemudian memantul sehingga ketinggiannya menjadi $\tfrac{1}{2}$
ketinggian sebelumnya.
Tentukan kebenaran setiap pernyataan berikut!

\par\smallskip
{\renewcommand{\arraystretch}{1.4}%
\begin{tabular}{@{}p{0.76\linewidth}cc@{}}
\hline
\textbf{Pernyataan} & \textbf{Benar} & \textbf{Salah}\\
\hline
(1)~Lintasan total sebelum pantulan pertama adalah 50 m.
  & $\bigcirc$ & $\bigcirc$\\
(2)~Panjang lintasan total tepat setelah pantulan ke-4 adalah
    137,5 m.
  & $\bigcirc$ & $\bigcirc$\\
(3)~Panjang lintasan total tepat setelah pantulan ke-3 sudah
    melebihi 137 m.
  & $\bigcirc$ & $\bigcirc$\\
(4)~Lintasan pertama kali melebihi 137 m tepat setelah
    pantulan ke-4.
  & $\bigcirc$ & $\bigcirc$\\
\hline
\end{tabular}}
\end{soal}

\begin{soal}{28}  % PG Majemuk
Sebuah balok kecil bermassa 102 gram dilepaskan dari titik A.
Balok bergerak melewati lintasan lengkung AB berbentuk seperempat
lingkaran jari-jari $R$, kemudian bergerak di lintasan datar BC
dengan $\mu_k=0{,}45$ dan tepat berhenti di C.
\gambar[0.55\textwidth]{images/1/1-11.png}
Tentukan kebenaran setiap pernyataan berikut!

\par\smallskip
{\renewcommand{\arraystretch}{1.4}%
\begin{tabular}{@{}p{0.76\linewidth}cc@{}}
\hline
\textbf{Pernyataan} & \textbf{Benar} & \textbf{Salah}\\
\hline
(1)~Energi potensial yang hilang pada lintasan AB adalah $mgR$.
  & $\bigcirc$ & $\bigcirc$\\
(2)~Kerja yang dilakukan gaya gesek pada BC adalah $-\mu_k mg\cdot BC$.
  & $\bigcirc$ & $\bigcirc$\\
(3)~Perbandingan $R/BC = \mu_s = 0{,}5$.
  & $\bigcirc$ & $\bigcirc$\\
(4)~Perbandingan $R/BC = \mu_k = 0{,}45$.
  & $\bigcirc$ & $\bigcirc$\\
\hline
\end{tabular}}
\end{soal}

\begin{soal}{29}  % PG Biasa
Pada $t=0$ s, ambulans bergerak dengan kecepatan konstan menuju
suatu perempatan yang berjarak 1 km di depannya sambil membunyikan
sirine 1000 Hz. Tepat saat 1 menit, frekuensi sirine yang
dideteksi oleh sensor di perempatan adalah \ldots\ Hz
(cepat rambat bunyi 340 m/s).
\begin{pilB}
  \item $951{,}5$ \item $1000{,}0$ \item $1025{,}0$ \item $1051{,}5$ \item $1100{,}0$
\end{pilB}
\end{soal}

\begin{soal}{30}  % PG Majemuk
Sebuah balok massa $m_1$ meluncur di meja licin, terhubung oleh
tali tak bermassa melalui katrol tak bermassa ke bola massa $m_2$
yang tergantung.
\gambar[0.5\textwidth]{images/1/1-12.png}
Tentukan kebenaran setiap pernyataan berikut!

\par\smallskip
{\renewcommand{\arraystretch}{1.4}%
\begin{tabular}{@{}p{0.76\linewidth}cc@{}}
\hline
\textbf{Pernyataan} & \textbf{Benar} & \textbf{Salah}\\
\hline
(1)~Percepatan sistem adalah $a = \dfrac{m_2 g}{m_1 + m_2}$.
  & $\bigcirc$ & $\bigcirc$\\
(2)~Besar gaya tegangan tali sama di kedua sisi katrol
    (karena tali tak bermassa).
  & $\bigcirc$ & $\bigcirc$\\
(3)~Gaya tegangan tali adalah $T = m_2 g$.
  & $\bigcirc$ & $\bigcirc$\\
(4)~Gaya tegangan tali adalah $T = \dfrac{m_1 m_2 g}{m_1 + m_2}$.
  & $\bigcirc$ & $\bigcirc$\\
\hline
\end{tabular}}
\end{soal}

\begin{soal}{31}  % PG Biasa
Sebuah pesawat pengebom terbang horizontal dengan laju konstan
$v_p=60$ m/s pada ketinggian $h=128$ m. Tepat di atas titik O,
sebuah bom dilepaskan dan mengenai truk yang bergerak dengan
kecepatan konstan. Saat bom dilepaskan, truk berada pada jarak
$x_0=100$ m dari O. Jika tinggi truk adalah 3 m, laju truk
adalah \ldots\ m/s.
\gambar[0.55\textwidth]{images/1/1-13.png}
\begin{pilB}
  \item $20$ \item $25$ \item $32$ \item $40$ \item $48$
\end{pilB}
\end{soal}

\begin{soal}{32}  % PG Biasa
Sebuah balok bermassa $m=2$ kg diberi gaya konstan $F=50$ N dengan
arah membentuk sudut $\theta$ ke bawah terhadap horizontal
$\left(\tan\theta=\tfrac{3}{4}\right)$. Koefisien gesek statis
dan kinetis masing-masing 0,4 dan 0,3.
\gambar[0.55\textwidth]{images/1/1-14.png}
Besar percepatan balok adalah \ldots\ m/s$^2$.
\begin{pilB}
  \item $8{,}0$ \item $10{,}0$ \item $11{,}0$ \item $12{,}5$ \item $15{,}0$
\end{pilB}
\end{soal}

\begin{soal}{33}  % Isian
Mobil pemadam kebakaran bergerak menuju perempatan jalan dengan
kelajuan 72 km/jam sambil menyalakan sirine. Sensor di perempatan
mendeteksi perubahan frekuensi sirine tepat sebelum dan setelah
mobil melewatinya sebesar 87 Hz. Cepat rambat bunyi di udara
340 m/s. Frekuensi sirine mobil pemadam tersebut adalah \ldots\ Hz.
\end{soal}

\begin{soal}{34}  % Isian
Sebuah benda bermassa 6 kg memiliki energi kinetik sebesar 200 J
sesaat sebelum menumbuk tanah. Benda tersebut dijatuhkan dari
ketinggian \ldots\ m. (Hambatan udara diabaikan.)
\end{soal}

\begin{soal}{35}  % PG Majemuk
Perhatikan sistem pada gambar. Massa benda A adalah 10 kg dan
massa benda B adalah 5 kg. Benda A ditarik ke atas sehingga
sistem bergerak ke atas dengan percepatan konstan 2 m/s$^2$.
\gambar[0.32\textwidth]{images/1/1-15.png}
Tentukan kebenaran setiap pernyataan berikut!

\par\smallskip
{\renewcommand{\arraystretch}{1.4}%
\begin{tabular}{@{}p{0.76\linewidth}cc@{}}
\hline
\textbf{Pernyataan} & \textbf{Benar} & \textbf{Salah}\\
\hline
(1)~Massa total sistem A dan B adalah 15 kg.
  & $\bigcirc$ & $\bigcirc$\\
(2)~Persamaan gerak sistem: $T_1 - Mg = Ma$, dengan $M=15$ kg.
  & $\bigcirc$ & $\bigcirc$\\
(3)~Besar tegangan tali $T_1 = 150$ N.
  & $\bigcirc$ & $\bigcirc$\\
(4)~Besar tegangan tali $T_1 = 180$ N.
  & $\bigcirc$ & $\bigcirc$\\
\hline
\end{tabular}}
\end{soal}

\begin{soal}{36}  % PG Majemuk
Sebuah gelombang transversal memiliki persamaan simpangan
\[
  y(t,x) = (0{,}23\,\text{m})\sin\!\left(\frac{\pi x}{5} - 8\pi t\right)
\]
dengan $t$ dalam sekon dan $x$ dalam meter. Tentukan kebenaran
setiap pernyataan berikut!

\par\smallskip
{\renewcommand{\arraystretch}{1.4}%
\begin{tabular}{@{}p{0.76\linewidth}cc@{}}
\hline
\textbf{Pernyataan} & \textbf{Benar} & \textbf{Salah}\\
\hline
(1)~Gelombang merambat ke arah sumbu $+x$.
  & $\bigcirc$ & $\bigcirc$\\
(2)~Amplitudo kecepatan transversal tali adalah
    $A\omega = 0{,}23(8\pi) \approx 5{,}78$ m/s.
  & $\bigcirc$ & $\bigcirc$\\
(3)~Besar laju transversal di $x=1{,}5$ m, $t=0{,}5$ s lebih besar
    dari 4 m/s.
  & $\bigcirc$ & $\bigcirc$\\
(4)~Besar laju transversal di $x=1{,}5$ m, $t=0{,}5$ s adalah
    $|0{,}23(8\pi)\cos(0{,}3\pi)|\approx 3{,}40$ m/s.
  & $\bigcirc$ & $\bigcirc$\\
\hline
\end{tabular}}
\end{soal}

\begin{soal}{37}  % PG Majemuk
Gelombang stasioner pada seutas tali dinyatakan oleh persamaan
\[
  y = 5\cos\!\left(\frac{\pi x}{3}\right)\sin(40\pi t)
\]
di mana $x$ dan $y$ dalam cm dan $t$ dalam detik. Tentukan
kebenaran setiap pernyataan berikut!

\par\smallskip
{\renewcommand{\arraystretch}{1.4}%
\begin{tabular}{@{}p{0.76\linewidth}cc@{}}
\hline
\textbf{Pernyataan} & \textbf{Benar} & \textbf{Salah}\\
\hline
(1)~Gelombang stasioner di atas terbentuk dari superposisi dua
    gelombang berjalan berlawanan arah.
  & $\bigcirc$ & $\bigcirc$\\
(2)~Faktor spasial gelombang stasioner ini adalah $\cos(\pi x/3)$.
  & $\bigcirc$ & $\bigcirc$\\
(3)~Jarak antara dua simpul berturut-turut adalah $\lambda/4$.
  & $\bigcirc$ & $\bigcirc$\\
(4)~Jarak antara dua simpul berturut-turut pada gelombang
    stasioner ini adalah 3 cm.
  & $\bigcirc$ & $\bigcirc$\\
\hline
\end{tabular}}
\end{soal}

\begin{soal}{38}  % Isian
Grafik di bawah menyatakan simpangan gelombang berjalan pada $t=0$
yang merambat ke kanan. Jika frekuensi gelombang adalah 5 Hz,
tuliskan persamaan simpangan gelombang tersebut!
\gambar[0.55\textwidth]{images/1/1-16.png}
$y(x,t) = $ \underline{\hspace{6cm}}\ (satuan SI)
\end{soal}

\begin{soal}{39}  % PG Biasa
Sebuah senar gitar nilon panjang 90 cm dijepit kedua ujungnya
sehingga menghasilkan gelombang berdiri dengan total 4 buah simpul.
Laju rambat gelombang transversal pada senar adalah 132 m/s.
Frekuensi dari gelombang berjalan penyusun gelombang berdiri
tersebut adalah \ldots\ Hz.
\begin{pilB}
  \item $110$ \item $165$ \item $220$ \item $330$ \item $440$
\end{pilB}
\end{soal}

\begin{soal}{40}  % Isian  [SOAL BARU: energi pegas]
Sebuah pegas dengan konstanta pegas $k=200$ N/m dimampatkan sejauh
$x=0{,}3$ m dari posisi setimbang kemudian dilepaskan untuk mendorong
sebuah bola bermassa $m=2$ kg di atas permukaan horizontal licin.
Tentukan kelajuan bola saat pegas kembali ke posisi setimbang
dalam satuan m/s.

\vspace{0.3cm}
$v = $ \underline{\hspace{4cm}}\ m/s
\end{soal}

\end{document}
